\subsection{Proof of Theorem \ref{thm2}}
We will show here that the expansion constructed in the previous section provides an approximate solution of the true solution $u^\eps$ of \eqref{limit_equation}. Given $m \in \N$, we~first introduce $\tilde{u}^{m+2}$ the solution in $\dot{H}^1(\Omega^\eps)$ of the Stokes problem
\begin{equation*}
\left\{
\begin{aligned}
-\Delta \tilde{u}^{m+2} + \na \tilde{p}^{m+2} & = 0 \quad \text{in } \Omega^\eps, \\
\div \tilde{u}^{m+2} & = - \pa_1 u^{m+2}_1 1_{\Omega^0} \quad \text{in } \Omega^\eps, \\
\tilde{u}^{m+2}\vert_{\pa \Omega^\eps} & = 0.
\end{aligned}
\right.
\end{equation*}
Such $\tilde{u}^{m+2}$ exists, \cf \cite[Th.\,IV.3.3]{Galdi}, and satisfies
\[
\|\na \tilde{u}^{m+2}\|_{L^2(\Omega^\eps)} \lesssim \|\pa_1 u^{m+2}_1 1_{\Omega^0} \|_{L^2(\Omega^\eps)} \lesssim 1.
\]
We then define $(u^\eps_{\app}, p^{\eps}_{\app})$ as follows:
\begin{equation*}
\begin{aligned}
(u^\eps_{\app},p^\eps_{\app})(x,y,z) & = \sum_{i=0}^{m+1} \eps^i (u^i, p^i)(x,y,z) + \eps^{m+2} (u^{m+2}_1,0,0,0)(x,y,z) \\
& \hspace*{1cm}+ \eps^{m+2} (\tilde{u}^{m+2},\tilde{p}^{m+2})(x,y,z), \qquad X = (x,y,z) \in \Omega^0, \\
(u^\eps_{\app}, p^\eps_{\app})(x,y,z) & = \sum_{i=0}^{m+1} \eps^i (U^i, P^i)(x/\eps,y,z) + \eps^{m+2} (U^{m+2}_1,0,0,0)(x/\eps,y,z) \\
& \hspace*{1cm}+ \eps^{m+2} (\tilde{u}^{m+2},\tilde{p}^{m+2})(x,y,z), \qquad \: X = (x,y,z) \in \Omega^\eps\setminus\Omega^0.
\end{aligned}
\end{equation*}
We then introduce $v^\eps = u^\eps - u^\eps_{\app}$, $q^\eps = p^\eps - p^\eps_{\app}$, solving
\begin{align}
-\Delta v^\eps + \na q^\eps & = - \eps^{m+2} (\Delta u^{m+2}_1,0,0) \quad \text{ in } \: \Omega^0,\notag \\
-\Delta v^\eps + \na q^\eps & = - \sum_{i=m}^{m+1} \eps^i (\pa^2_2 + \pa^2_3) (\eps U^{i+1}_1,U^i_2,U^i_3) - \sum_{i=m}^{m+1} \eps^i (0, \pa_2 P^i, \pa_3 P^i) \quad \text{ in } \: D^\eps, \notag \\
\div v^\eps & = 0 \quad \text{ in } \: \Omega^\eps, \\
[v^\eps]\vert_{\pa \Omega^0} & = 0, \notag\\
[\Sigma(v^\eps,q^\eps)n]\vert_{\pa \Omega^0} & = - \eps^{m+1}(\eps \pa_1 u^{m+2}_1\vert_{x_1=0},\pa_1 u^{m+1}_2\vert_{x_1=0}, \pa_1 u^{m+1}_3\vert_{x_1=0}), \notag \\
v^\eps \vert_{\pa \Omega^\eps} & = 0. \notag
\end{align}
Performing a simple energy estimate, we~find
\begin{align*}
\int_{\Omega^\eps} |\na v^\eps|^2 = \eps^{m+2} \int_{\Omega^0} \na u_1^{m+2} : \na v_1^\eps - \int_{D^\eps} R^\eps \cdot v^\eps + \int_{\pa \Omega^0} r^\eps \cdot v^\eps,
\end{align*}
where
\begin{align*}
R^\eps & = - \sum_{i=m}^{m+1} \eps^i (\pa^2_2 + \pa^2_3) (\eps U^{i+1}_1,U^i_2,U^i_3) - \sum_{i=m}^{m+1} \eps^i (0, \pa_2 P^i, \pa_3 P^i) = O(\eps^m), \\
r^\eps & = \eps^{m+1}(\eps \pa_1 u^{m+2}_1\vert_{x_1=0},\pa_1 u^{m+1}_2\vert_{x_1=0}, \pa_1 u^{m+1}_3\vert_{x_1=0}) = O(\eps^{m+1}).
\end{align*}
By the Poincaré and trace inequalities
\[
\| v^\eps \|_{L^2(D^\eps)} \lesssim \eps \|\na v^\eps \|_{L^2(D^\eps)}, \quad \| v^\eps \|_{L^2(\pa\Omega^0)} \lesssim \sqrt{\eps} \|\na v^\eps \|_{L^2(D^\eps)},
\]
we infer easily
\[
\|\na v^\eps \|_{L^2(\Omega^\eps)} \lesssim \eps^{m+1}.
\]
Restricting to $\Omega^0$, Theorem \ref{thm2} follows.

